Denne masteroppgaven tar for seg to hovedspørsmål: hva er den forventede levetiden til et sammarbeidende GNSS-R-formasjonsflygingsoppdrag med to satellitter, og hvordan bør denne estimerte  levetiden tolkes sett i lys av prediksjonsnøyaktigheten til det underliggende Basilisk-rammeverket? For å besvare disse spørsmålene ble det utviklet to separate simulatorer: en simulator for å sammenlikne propagasjonssimuleringer, og en GNSS-R-formasjonsflygingssimulator. For å estimere levetiden til formasjonsoppdraget måtte begrensningene tilknyttet GNSS-R oppdraget modelleres, inkludert satellittkonfigurasjon, det elektriske systemet, autonom modusbytting, attityde- og formasjonskontroll. Den estimerte levetiden ble predikert ved hjelp av modeller utledet fra to simuleringsrunder, som henholdsvis undersøkte drivstofforbrukets sensitivitet for leder- og følgersatellittenes utskytningshastighet fra et delt romfartøy, og drivstofforbruket knyttet til formasjonsvedlikehold i stasjonær tilstand. De endelige levetidsestimatene, utledet ved bruk av disse modellene, var $413.6~\mathrm{dager}$ og $366.5~\mathrm{dager}$ for henholdsvis det mest optimistiske og det mest pessimistiske utskytningsscenarioet. På grunn av modelleringsusikkerhet og observert formasjonsdrift ble estimatet vurdert som konservativt.
 
Basilisk-propagasjonskonfigurasjonen som ble brukt i levetidsanalysen, ble utviklet som en videreføring av spesialiseringsprosjektet, hvor den tidligere eksponentielle atmosfæremodellen ble erstattet av NRLMSISE-00  atmosfæremodellen. Denne konfigurasjonen ble sammenlignet med en kontinuerlig oppdatert Skyfield-\gls{sgp4}-referanse ved å evaluere både tilstandsforskjeller for samme satellitt og formasjonsdriftsforskjeller. Sammenligningen av samme satellitt viste at konfigurasjonen brukt i denne masteroppgaven divergerte betydelig mindre fra Skyfield-\gls{sgp4}-referansen enn konfigurasjonen fra spesialiseringsprosjektet. Sammenligningen av formasjonsdrift viste imidlertid at konfigurasjonen i denne masteroppgaven predikerte større passiv leder-følger-drift enn Skyfield-\gls{sgp4}-referansen. Siden levetidskampanjene dermed sannsynligvis var utsatt for større relativ drift enn det som ble predikert av sammenligningsreferansen, støtter dette vurderingen av levetidsestimatet som konservativt.
